\subsection{型关于正埃尔米特型的约化。}

在本节中，我们将假设 E 在 A 上的维数为有限数 n，并以 Φ 表示 E 上的一个非退化正埃尔米特型。

由于从 E 到 \(\mathbf{E}^*\) 中的与 \(\Phi\) 相伴的线性映射都是双射，故对 E 中任意元素 \(x, y, z\) （ \(y \neq 0\) ），存在 E 的恰好一个元素 \(t\) 使 \(\Phi(x, z) = \overline{\Phi(t, y)}\) 。特别地，若 \(\mathrm{A} = \mathrm{K}\) 或 \(\mathrm{A} = \mathrm{K}(i)\) ，且 \(u\) 是 E 到自身的（关于 J 的）半线性映射（第 II 章，附录 I，第 1 小节），则对每个 \(x \in \mathrm{E}\) ，存在恰好一个元素 \(u^*(x)\) ，使得对一切 \(y \in \mathrm{E}\) 有 \(\Phi(x, u(y)) = \overline{\Phi(u^*(x), y)}\) 。立即可以看出，\(u^*\) 是 E 到自身的（关于 J 的）半线性映射；称之为 \(u\) 的伴随。

注记. --- 1) 当 J 是恒等映射时，我们重新得到 § 1 第 8 小节中定义的同态的伴随概念。

\begin{enumerate}
\def\labelenumi{\arabic{enumi})}
\setcounter{enumi}{1}
\tightlist
\item
  仍假设 A = K 或 A = K(i)。设 E \(^{j}\) 是 § 1 第 2 小节定义 5 中定义的右 A-模。型 Φ 是 E × E \(^{j}\) 上的双线性型，Φ \(^{j}\) 是 E \(^{j}\) × E 上的双线性型，而 u 是从 E 到 E \(^{j}\) 中的 A-线性映射。这个同态在 § 1 第 8 小节意义下的伴随是从 E 到 E \(^{j}\) 中的一个 A-线性映射；我们立即看出它与上面定义的映射 u* 重合。
\end{enumerate}

我们称 E 的自同态 \(u\) 是（关于 \(\Phi\) 的）正规自同态，如果有 \(uu^{*}=u^{*}u\) 。
% semantic-fragments: legacy-input-seam
\relax{}
\paragraph{正规自同态的例：}

\begin{enumerate}
\def\labelenumi{\arabic{enumi})}
\item
  关于 \(\Phi\) 的酉自同构（\(\S 6\)，第 2 小节），它们由关系 \(u^{-1} = u^{*}\) 刻画（\(\S 1\)，第 8 小节，命题 8 的推论）；
\item
  自同态 \(u\)，满足 \(u^{*}=u\)；它们称为埃尔米特自同态。
\end{enumerate}

对 E 的每个埃尔米特自同态 \(u\)，令 \(\Phi_{u}(x,y)=\Phi(u(x),y)\)；我们有

\[
\Phi_ {u} (y, x) = \Phi (u (y), x) = \Phi (y, u (x)) = \overline {{\Phi (u (x) , y)}} = \overline {{\Phi_ {u} (x , y)}}.
\]

这表明 \(\Phi_{u}\) 是 E 上的埃尔米特型。反之，设 \(\Psi\) 是 E 上的一个埃尔米特型；由于从 E 到 E* 的映射 \(s_{\Phi}\)（它与 \(\Phi\) 相伴）是双射，故对每个 \(x \in \mathbf{E}\) ，存在 E 中唯一的元素 \(u(x)\) 使得 \(\Psi(x, y) = \Phi(u(x), y)\)；容易验证 \(u\) 是 E 的埃尔米特自同态。于是 \(u \to \Phi_{u}\) 是从 E 的埃尔米特自同态所成的集到 E 上埃尔米特型所成的集上的双射，称为典范双射。

设 A = K 或 A = K(i)；给定 E 上的一个双线性型 Ψ，对每个 \(x \in E\) ，存在唯一的元素 \(u(x)\)，即 E 中唯一的元素，使得 Ψ(x, y) = \(\overline{\Phi(u(x), y)}\)；立即可以看出 u 是 E 到其自身的（关于 J 的）半线性映射。由于 \(\overline{\Phi(u(x), y)} = \Phi(y, u(x)) = \overline{\Phi(u^{*}(y), x)}\)，可以看出，要使 Ψ 对称（相应地，交错），必须且只需 \(u^{*} = u\)（相应地，\(u^{*} = -u\)）。

定理 2. --- 设 S 是 E 的一族自同态（相应地，当 A = K 或 A = K(i) 时，为 E 到 E 中的半线性映射所成的集），它在映射 u → u* 下稳定。那么，若 V 是 E 的一个在 S 下稳定的子空间，则其正交补 V⁰ 也在 S 下稳定。此外，E 是在 S 下稳定、在 ≠ \{0\} 且在 S 下稳定的诸子空间中为极小的、且两两正交的一些子空间的直和。

实际上，设 V 是 E 的一个在 S 下稳定的子空间；对任意 \(x \in \mathrm{V}^0\) 、\(y \in \mathrm{V}\) 及 \(u \in \mathrm{S}\)，有 \(u^*(y) \in \mathrm{V}\)，从而 \(\Phi(y, u(x)) = \Phi(u^*(y), x) = 0\)（相应地，\(\Phi(y, u(x)) = \overline{\Phi(u^*(y), x)} = 0\)），于是 \(u(x) \in \mathrm{V}^0\)；因而 \(\mathrm{V}^0\) 在 S 下稳定，这就证明了第一个断言。对第二个断言，我们对 E 的维数 \(n\) 作归纳，\(n = 0\) 的情形是平凡的。当 \(n \neq 0\) 时，存在 E 的一个子空间 \(\mathrm{V} \neq \{0\}\)，它在 S 下稳定且极小，例如一个 \(\neq \{0\}\) 的、维数极小的稳定子空间。这时只需对 \(\mathrm{V}^0\) 应用归纳假设，因为 \(u\) 在 \(\mathrm{V}^0\) 上的限制（关于 \(\Phi\) 的限制）的伴随，等同于 \(\mathrm{V}^0\) 上 \(u\) 的伴随的限制。

推论 1. --- 设 A = K 或 A = K(i)。设 B 是 \(\mathfrak{L}_{\mathbf{A}}(\mathbf{E})\) 的一个在映射 \(u \to u^{*}\) 下稳定的子代数。则 E 是半单 B-模，且是两两正交的一些单子模的直和。代数 B 是半单的。

实际上，由于 E 的每个子 B-模 V 都有一个补，例如 \(V^{0}\)，故 B-模 V 是半单的（第八章 § 3 第 3 小节命题 7）。由于 E 的每个 \(\neq\{0\}\) 且极小的子 B-模都是单的，故 E 是两两正交的一些单子模的直和。最后，B 是半单代数，因为它有一个半单且忠实的模 E，其反向模是有限生成的（第八章 § 5 第 1 小节命题 3）。

推论 2. --- 在推论 1 的假设与记号下，再设代数 B 是交换的。则其所有元素都是 E 的（正规）半单自同态。当 A = K（相应地，A = K(i)）时，E 是两两正交的一些单 B-子模的直和，这些子模是 A 上维数为 1 或 2（相应地，1）的向量空间。

第一个断言由第八章 § 9 第 1 小节命题 2 得出，因为 B 是半单的。另一方面，每个单 B-模同构于形如 B/m 的 B-模，其中 m 是 B 的一个极大理想，从而同构于 A 上有限次的一个域 L；当 A = K（相应地，A = K(i)）时，域 L 必然同构于 K 或 K(i)（相应地，K(i)），因为 K(i) 是代数闭的（第六章 § 2 第 6 小节定理 3）。

命题 4. --- 设 u 是 E 的一个正规自同态。当 A 等于 K(i) 或 K 上的四元数体时，存在由 u 的特征向量组成的 E 的（关于 Φ 的）标准正交基。当 A = K 时，u 是半单的，且 E 是在 u 下稳定、两两正交、维数为 1 或 2 的一些子空间的直和。

先考察 A 交换的情形（A = K 或 A = K(i)）。这时子代数 B = A{[}u, u*{]}（ℒ\_A(E) 的）是交换的，因为 u 正规；由 § 1 第 8 小节的公式 (32) 与 (33)，它在映射 v → v* 下稳定。关于 A = K 情形的断言随即由定理 2 的推论 2 得出。当 A = K(i) 时，该推论还表明 E 是两两正交、在 u 下稳定、维数为 1 的一些向量子空间 Ax\_i (i = 1, . . ., n) 的直和；若令 \(e_{i} = (\Phi(x_{i}, x_{i}))^{-1/2}x_{i}\)，则 \((e_{i})\) 就是所要的标准正交基。

当 A 是 K 上的四元数体时，同样只需依定理 2 证明：在 u 与 u* 下稳定的、≠ \{0\} 的子空间所成的集的每个极小元素都具有维数 1。而今，这样的子空间 V 必含 u 的一个特征向量 x ≠ 0 (*)，这可如下看出：注意四元数体 A 含 K(i) 作为代数闭子域，并把 V 的标量域限制到 K(i)。于是只需证明 \(x\)（\(u\) 的特征向量）也是 \(u^{*}\) 的特征向量。设 \(u(x) = ax (a \in A)\)；这时我们有 \(\Phi(u(x), x) = a\Phi(x, x) = \Phi(x, x)a = \Phi(x, \bar{a}x)\)，因为 \(\Phi(x, x)\) 属于 A 的中心；而另一方面 \(\Phi(u(x), x) = \Phi(x, u^*(x))\)；由此得 \(\Phi(x, u^*(x) - \bar{a}x) = 0\)，因此可以写 \(u^*(x) = \bar{a}x + z\)，其中 \(z\) 是与 \(x\) 正交的一个向量。于是我们有

\[
\Phi (u ^ {*} (x), u ^ {*} (x)) = \bar {a} a \Phi (x, x) + \Phi (z, z) = \Phi (u (x), u (x)) + \Phi (z, z).
\]

而今，由于 \(u\) 是正规的，我们有

\[
\begin{array}{r l} \Phi (u (x), u (x)) & = \Phi (x, u ^ {*} u (x)) = \Phi (x, u u ^ {*} (x)) \\ & = \Phi (u u ^ {*} (x), x) = \Phi (u ^ {*} (x), u ^ {*} (x)). \end{array}
\]

于是有 \(\Phi(z, z) = 0\)，按假设这蕴含 \(z = 0\) 及 \(u^{*}(x) = \bar{a}x\)。证毕。

注记. --- 由定理 2 可知，对应于 \(u\) 的两个不同特征值的特征子空间是正交的。

命题 5. --- 设 u 是 E 的埃尔米特自同态。u 的特征值都属于 K，且存在由 u 的特征向量组成的 E 的标准正交基。

当 A 等于 K(i) 或 K 上的四元数体时，依命题 4，只需证明第一个断言；而今，若 x 是 u 的一个非零特征向量，a 是相应的特征值，则由假设 u = u*，我们有

\[
a \Phi (x, x) = \Phi (u (x), x) = \Phi (x, u (x)) = \Phi (x, x) \bar {a};
\]

由于 \(\Phi(x, x)\) 是 A 的中心中的一个非零元素，由此得 \(a = \bar{a}\)，从而 \(a \in K\)。因此，埃尔米特矩阵的所有特征值都在 K 中。现在设 A 等于 K。矩阵 \(M\)（即 \(u\) 关于 E 的一个标准正交基的矩阵）这时是对称的；于是，若把它看作 \(K(i)\) 上的矩阵，它就是一个埃尔米特矩阵。证明的第一部分于是表明，这个矩阵的所有特征值都在 K 中，从而当 \(E \neq \{0\}\) 时，它在 E 中有 \(\neq 0\) 的特征向量。由此推出，每个在 \(u\) 下稳定且极小的子空间必然具有维数 1，而结论如命题 4 中那样立即可得。

命题 6. --- a) 设 Ψ 是 E 上的一个埃尔米特型（相应地，当 A = K 或 A = K(i) 时，为对称双线性型）。存在 E 的一个基，它关于 Φ 是标准正交的，关于 Ψ 是正交的。

\begin{enumerate}
\def\labelenumi{\alph{enumi})}
\setcounter{enumi}{1}
\tightlist
\item
  设 A = K 或 A = K(i)，并设 Ψ 是 E 上的一个交错双线性型。存在 E 的一个基，它关于 Φ 是标准正交的，且 Ψ 关于这个基的矩阵具有以下形状
\end{enumerate}

\[
\left( \begin{array}{c c c c} 0 & a _ {1} & 0 & 0 \ldots 0 \\ - a _ {1} & 0 & 0 & 0 \ldots 0 \\ 0 & 0 & 0 & a _ {2} \ldots 0 \\ 0 & 0 & - a _ {2} & 0 \ldots 0 \\ \ldots \ldots \end{array} \right).
\]

其中诸 \(a_{i}\) 在 K 中 \(\geqslant 0\)。

当 \(\Psi\) 是埃尔米特型时，我们的断言由命题 5 以及 E 上埃尔米特型与关于 \(\Phi\) 的埃尔米特自同态之间的典范对应立即得出。在另外两种情形，设 \(u\) 是本小节开头由公式 \(\Psi(x, y) = \overline{\Phi(u(x), y)}\) 所定义的 E 到其自身的半线性映射。于是 \(u^2\) 是一个 A-线性映射；当 \(\Psi\) 对称（相应地，交错）时，有 \(u = u^*\)（相应地，\(u = -u^*\)），因而 \(u^2\) 是埃尔米特的；从而也有

\[
\begin{array}{r l} \Phi (u ^ {2} (x), x) & = \overline {{\Phi (x , u ^ {2} (x))}} = \Phi (u ^ {*} (x), u (x)) = \Phi (u (x), u (x)) \\ & (\text { resp. } - \Phi (u (x), u (x))) \end{array}
\]

这表明埃尔米特型 \((x, y) \to \Phi(u^{2}(x), y)\) 是正的（相应地，负的）。然后应用定理 2，设 V 是 E 的 \(\neq\{0\}\) 的、在 u 与 \(u^{*}\) 下稳定的子空间所成的族中的一个极小元素。由于 \(u^{2}\) 是埃尔米特自同态，V 含有一个 \(x \neq 0\)，它是 \(u^{2}\) 的特征向量；设 \(u^{2}(x) = ax\)，其中 \(a \in K\)（命题 5）。

当 \(\Psi\) 对称时，不等式 \(\Phi(u^{2}(x), x) \geqslant 0\) 表明 \(a \geqslant 0\)。设 \(y = a^{1/2}x + u(x)\)；我们有 \(u(y) = a^{1/2}u(x) + ax = a^{1/2}y\)，且 Ay 在 \(u\) 与 \(u^{*} = u\) 下稳定。若 \(y = 0\)，则 Ax 在 \(u\) 与 \(u^{*}\) 下稳定；在所有情形，V 都是维数为 1 的子空间，而我们的结论由定理 2 立即可得。

当 \(\Psi\) 交错时，有 \(a \leqslant 0\)；设

\[
y = (- a) ^ {1 / 2} x + u (x), \quad z = (- a) ^ {1 / 2} x - u (x).
\]

\[
\text { One has } u (y) = - (- a) ^ {1 / 2} z, u (z) = (- a) ^ {1 / 2} y \text { and }
\]

\[
\Phi (y, z) = - a \Phi (x, x) - \Phi (u (x), u (x)) = - \Phi (a x, x) + \Phi (u ^ {2} (x), x) = 0.
\]

若 \(y = z = 0\)，则有 \(a = 0\) 且 \(u(x) = 0\)，于是 \(Ax\) 在 \(u\) 与 \(u^{*}\) 下稳定，\(V\) 的维数为 1，且 \(\Psi\) 在 \(V\) 上的限制的矩阵为零，因为 \(\Psi\) 是交错的。否则，\(y\) 与 \(z\) 都 \(\neq 0\)，它们生成 \(V\)，且由于它们正交，\(V\) 的维数为 2，而 \(\Psi\) 在 \(V\) 上的限制的矩阵具有形状 \(\begin{pmatrix} 0 & b \\ -b & 0 \end{pmatrix}\)。最后，依公式 \(\Psi(x', y') = \overline{\Phi(u(x'), y')}\)，有 \(\Psi(x', y') = 0\)，只要 \(x'\) 与 \(y'\) 属于两个在 \(u\) 下稳定且关于 \(\Phi\) 正交的子空间。这就证明了存在 \(E\) 的（关于 \(\Phi\) 的）标准正交基，\(\Psi\) 关于它的矩阵具有所指出的形状。证毕。

习题. --- 1) 假设 § 7 开头的诸假设成立；设 \(\Phi\) 是 E 上的一个正埃尔米特型。

\begin{enumerate}
\def\labelenumi{\alph{enumi})}
\item
  要使 \(\Phi(x,x)\Phi(y,y)=\Phi(x,y)\overline{\Phi(x,y)}\)，必须且只需 \(x\) 与 \(y\) 线性相关，或由 \(x\) 与 \(y\) 生成的平面是迷向的。
\item
  假设对一切 \(x \in E\)，\(\Phi(x, x)\) 是 K 中的一个平方。证明对 E 中任意两个向量 x、y，我们有
\end{enumerate}

\[
\sqrt {\Phi (x + y , x + y)} \leqslant \sqrt {\Phi (x , x)} + \sqrt {\Phi (y , y)}.
\]

若 \(\Phi\) 非退化，则这个不等式的两边仅当 \(\alpha x + \beta y = 0\) 时才能相等，其中 \(\alpha\) 与 \(\beta\) 是 K 中两个不全为零的元素，且满足 \(\alpha \beta \leqslant 0\)。

\begin{enumerate}
\def\labelenumi{\arabic{enumi})}
\setcounter{enumi}{1}
\tightlist
\item
  假设 \(\mathbf{A} = \mathbf{K}\) 或 \(\mathbf{A} = \mathbf{K}(i)\)。设 \(X\) 是一个 \(n\) 阶埃尔米特方阵，元素在 \(\mathbf{A}\) 中，使得对每个非空子集 \(\mathbf{H}\) ⊂ \([1, n]\)，有 \(X_{\mathbf{H},\mathbf{H}} \geqslant 0\)（第 1 小节命题 3 的记号）。
\end{enumerate}

\begin{enumerate}
\def\labelenumi{\alph{enumi})}
\item
  设 \(\lambda\) 是 K 中一个 \(>0\) 的元素；证明矩阵 \(X + \lambda I\) 是正且非退化的（利用第 1 小节命题 3，对 \(n\) 作归纳论证）。
\item
  由此推出埃尔米特矩阵 X 是正的。
\end{enumerate}

\begin{enumerate}
\def\labelenumi{\arabic{enumi})}
\setcounter{enumi}{2}
\item
  假设 \(\mathbf{A} = \mathbf{K}\) 或 \(\mathbf{A} = \mathbf{K}(i)\)，且 E 是有限维的。设 \(\Phi_1, \Phi_2\) 是 E 上两个正埃尔米特型，\(V = (\alpha_{ij})\)、\(W = (\beta_{ij})\) 是这两个型关于 E 的同一基 \((e_i)\) 的矩阵。证明：埃尔米特型 \(\Phi\)（其矩阵 \((\gamma_{ij})\) 关于 \((e_i)\) 满足对每对指标 \(\gamma_{ij} = \alpha_{ij}\beta_{ij}\)）是正的；此外，若 \(\Phi_1\) 与 \(\Phi_2\) 非退化，则 \(\Phi\) 亦然。（在计算 \(\Phi(x, x)\) 时，把诸 \(\alpha_{ij}\) 用 \(\Phi_1(c_i, c_i)\) 的值表出，其中 \((c_i)\) 是 E 关于 \(\Phi_1\) 的一个正交基。）
\item
  设 A = K 或 A = K(i)。设 R 是 A 上一个 n 阶埃尔米特矩阵；称 R 具有符号差 \((s, t)\)，如果 \(A^{n}\) 上以 R 为其关于典范基的矩阵的埃尔米特型具有符号差 \((s, t)\)。以 \(\Delta_{k}\) 记从 R 中删去指标 \textgreater{} k 的行与列所得的 R 的主子式；假设 \(\Delta_{s+t} \neq 0\)，且对任何指标 \(k < s + t\)，\(\Delta_{k}\) 与 \(\Delta_{k+1}\) 不同时为零（参见 § 6 习题 1 b)）。证明：若 \(\Delta_{k} = 0\) 对某个 \(k < s + t\) 成立，则 \(\Delta_{k-1}\) 与 \(\Delta_{k+1}\) 符号相反（用 § 6 习题 1 a) 的方法），且数 s - t 等于
\end{enumerate}

\[
\operatorname{sgn} \Delta_ {1} + \operatorname{sgn} \left(\Delta_ {1} \Delta_ {2}\right) + \dots + \operatorname{sgn} \left(\Delta_ {s + t - 1} \Delta_ {s + t}\right)
\]

（利用 § 6 的习题 2）。

\begin{enumerate}
\def\labelenumi{\arabic{enumi})}
\setcounter{enumi}{4}
\item
  假设 A = K 或 A = K(i)；设 R、S 是 A 上两个埃尔米特矩阵，其符号差分别为 \((s, t)\) 与 \((s', t')\)（习题 4）；证明矩阵 \(R \otimes S\) 是符号差为 \((ss' + tt', st' + s't)\) 的埃尔米特矩阵。
\item
  设 K 是一个极大有序域，L 是 K 上一个有限维单代数。若 \((s, t)\) 是 L 上对称双线性型 \((x, y) \to \mathrm{Tr}_{\mathrm{L/K}}(xy)\) 的符号差，证明我们有：
\end{enumerate}

\(s-t=m\)，若 L 同构于 K 上一个 \(m\) 阶矩阵代数；

s - t = 0，若 L 同构于域 K(i) 上的矩阵代数；

\(s - t = -2m\)，若 L 同构于 K 上四元数体上的一个 \(m\) 阶矩阵代数。

\begin{enumerate}
\def\labelenumi{\arabic{enumi})}
\setcounter{enumi}{6}
\tightlist
\item
  假设 A 满足 § 7 开头的条件，且 E 具有有限维数 n。设 Φ 是 E 上的一个正且非退化的埃尔米特型。
\end{enumerate}

\begin{enumerate}
\def\labelenumi{\alph{enumi})}
\item
  证明每个关于 \(\Phi\) 的相似变换都可以以唯一的方式写成一个 K 中比 \(\geqslant 0\) 的位似与一个酉变换之积。
\item
  对 E 的每个基 \((a_{i})_{1\leqslant i\leqslant n}\)，证明存在一个且只有一个 E 的标准正交基 \((e_{i})_{1\leqslant i\leqslant n}\) 满足下列条件：\(1^{\circ}\) 对每个 \(m\)（\(1\leqslant m\leqslant n\)），由 \(a_1,\ldots ,a_m\) 生成的子空间等同于由 \(e_1,\ldots ,e_m\) 生成的子空间；\(2^{\circ}\) 在 K 中有 \(\Phi (a_i,e_i) > 0\)（对每个指标 \(i\)）（参见 § 6 第 1 小节命题 1）。
\item
  由 \(b\) 推出：对 A 上每个可逆方阵 \(M\)（\(n\) 阶），存在一个且只有一个矩阵对 \((L, U)\)（\(n\) 阶），使得 \(U\) 是酉矩阵，\(L = (\lambda_{ij})\) 的对角线下方全为零，且对角元素 \(\lambda_{ii}\) 属于 K 并 \(>0\)，最后 \(M = LU\)。
\end{enumerate}

设 A = K；设 L 是 Λ 上一个有限维埃尔米特空间，其度量型是正且非退化的。设 M、N 是 L 的两个子集，u 是从 M 到 N 上的一个双射，使得对 M 中任意点 a、b 有 \(e(u(a), u(b)) = e(a, b)\)，证明存在一个位移，它在 M 上的限制等于 u（按格拉姆–施密特正交化方法的方式，对由 M 生成的仿射线性簇的维数作归纳论证）。这个命题能否推广到 A 等于 K(i) 或 K 上四元数体的情形？

\begin{enumerate}
\def\labelenumi{\arabic{enumi})}
\setcounter{enumi}{8}
\item
  假设第 3 小节的诸条件成立，此外 A 是 K 上的四元数体。设 u 是 E 的一个正规自同态。证明 E 是两两正交的子空间 \(F_{k}\)（\(1 \leqslant k \leqslant r\)）的直和，使得在每个 \(F_{k}\) 中存在一个标准正交基（\(e_{ik}\)）（\(1 \leqslant i \leqslant n_{k}\)），且 \(u(e_{ik}) = \lambda_{k}e_{ik}\)（\(1 \leqslant i \leqslant n_{k}\)），而指标不同的两个 \(\lambda_{k}\) 不能通过 A 的一个内自同构互相变换。此外，若（\(F_{k}'\)）是 E 的一个具有同样性质的第二个分解，带有特征值组（\(\lambda_{k}'\)），则有 \(F_{k}' = F_{k}\)（在指标的一个置换下），且 \(\lambda_{k}' = \alpha_{k}\lambda_{k}\alpha_{k}^{-1}\)；u 关于特征值 \(\lambda_{k}\) 的特征向量所成的集合是 \(A_{k}\)（\(\lambda_{k}\) 在 A 中的中心化子）上由诸 \(e_{ik}\)（\(1 \leqslant i \leqslant n_{k}\)）生成的子空间。
\item
  假设第 3 小节的诸条件成立，此外 A 等于 K(i) 或 K 上的四元数体。设 u 是 E 的一个正规自同态。
\end{enumerate}

\begin{enumerate}
\def\labelenumi{\alph{enumi})}
\item
  证明：要使 E 的一个向量子空间 F 满足 \(u(\mathrm{F}) \subset \mathrm{F}\)，必须且只需 F 由 u 的特征向量生成；此时我们有 \(u(\mathrm{F}^{0}) \subset \mathrm{F}^{0}\) 和 \(u^{*}(\mathrm{F}) \subset \mathrm{F}\)。
\item
  要使 \(u\) 是酉的（相应地，\(u^{*} = u\)、\(u^{*} = -u\)），必须且只需对每个特征值 \(\lambda\)（\(u\) 的特征值）有 \(\lambda \bar{\lambda} = 1\)（相应地，\(\bar{\lambda} = \lambda\)、\(\bar{\lambda} = -\lambda\)）。
\item
  我们称 E 的埃尔米特自同态 \(u\) 是正的（相应地，正且非退化的），如果典范地对应于它的埃尔米特型 \((x, y) \to \Phi(u(x), y)\) 是正的（相应地，正且非退化的）；为此必须且只需 \(u\) 的所有特征值都 \(\geqslant 0\)（相应地，\(>0\)）。
\item
  证明对每个整数 m\textgreater{} 0，存在 E 的一个正规自同态 v 使得 \(v^{m}=u\)。若 u 是正埃尔米特的，则存在唯一的正埃尔米特自同态 v 使得 \(v^{m}=u\)，且存在一个多项式 \(f\in K[X]\) 使得 \(v=f(u)\)；这最后一个结果当 A = K 时也成立。
\end{enumerate}

\begin{enumerate}
\def\labelenumi{\arabic{enumi})}
\setcounter{enumi}{10}
\tightlist
\item
  假设第 3 小节的诸条件成立，此外 A = K。设 u 是 E 的一个正规自同态。
\end{enumerate}

\begin{enumerate}
\def\labelenumi{\alph{enumi})}
\item
  设 V 是 E 的不退缩为 0 且在 u 与 \(u^{*}\) 下稳定的子空间所成之集的一个极小元素；证明若 V 的维数为 2，则 u 在 V 上的限制是一个乘子 \textgreater{} 0 的直接相似。
\item
  证明 E 的每个在 u 下稳定的子空间也在 \(u^{*}\) 下稳定。（设 \(E_{0}\) 是把 E 的纯量域扩张到 \(\mathrm{K}(i)\) 后得到的向量空间；\(\Phi\) 是 \(E_{0}\) 上一个正非退化埃尔米特型在 E 上的限制，而 u 是 \(u_{0}\)（\(E_{0}\) 的一个对这个型而言的正规自同态）在 E 上的限制；此外，E 是满足下列条件的 \(x \in E_{0}\) 所成的集合：x 在 \(E_{0}\) 的一个对合半线性双射 j 下不变，且我们有 \(u_{0}j = ju_{0}\)。然后应用习题 10 a)。）
\end{enumerate}

\begin{enumerate}
\def\labelenumi{\arabic{enumi})}
\setcounter{enumi}{11}
\item
  假设第 3 小节的诸条件成立，此外 A 等于 K(i) 或 K 上的四元数体。证明：对 E 的每个自同态 u，存在 E 的一个标准正交基，使 u 关于这个基的矩阵在对角线下方只有零。（对 E 的维数作归纳论证，考虑 u 的一个特征向量。）
\item
  设 A 是一个满足 § 7 开头条件的体，E、F 是 A 上两个有限维向量空间，\(\Phi\)（相应地，\(\Psi\)）是 E（相应地，F）上的一个正非退化埃尔米特型。证明：对每个从 E 到 F 的线性映射 \(u\)，其伴随 \(u^{*}\)（\(u\) 的伴随，对 \(\Phi\) 与 \(\Psi\) 而言；参见 § 1 第 8 小节）具有如下性质：\(u^{*}u\) 与 \(uu^{*}\) 分别是 E 与 F 的正埃尔米特自同态（习题 10 c)）。
\item
  假设第 3 小节的诸条件成立。设 u 是 E 的一个自同态，\(h_{1}\)、\(h_{2}\) 是 E 的两个正埃尔米特自同态，使得我们有 \(h_{1}^{2}=u^{*}u\)、\(h_{2}^{2}=uu^{*}\)（习题 13 与 10 d)）。
\end{enumerate}

\begin{enumerate}
\def\labelenumi{\alph{enumi})}
\item
  证明存在一个酉自同态 \(\nu\) 使得 \(u = \nu h_{1} = h_{2}\nu\)，特别地 \(h_{1}\) 与 \(h_{2}\) 相似（注意 \(\bar{u}(0) = \bar{h}_{1}^{-1}(0)\)，且若 V 是与 \(\bar{u}(0)\) 正交的子空间，则 \(\Phi(u(x), u(x)) = \Phi(h_{1}(x), h_{1}(x))\) 对一切 \(x \in \mathrm{V}\) 成立。要使 \(\nu\) 被唯一确定，必须且只需 \(u\) 是双射的。要能取 \(\nu\) 与 \(h_{1}\) 交换，必须且只需 \(u\) 是正规的。
\item
  由 a) 推出：A 上每个 n 阶方阵 M 都可以写成 UDV，其中 U 与 V 是酉矩阵，D 是一个对角矩阵，其元素 \(\geqslant 0\) 且以 MM* 的特征值作为其平方。
\end{enumerate}

\begin{enumerate}
\def\labelenumi{\arabic{enumi})}
\setcounter{enumi}{14}
\tightlist
\item
  假设第 3 小节的诸条件成立。证明 A 上每个正埃尔米特矩阵 H 都可以写成 \(LL^{*}\) 的形式，其中 \(L = (\lambda_{ij})\) 的对角线下方只有零，且对角线项属于 K 并 \(\geqslant 0\)；此外，当 H 可逆时 L 由这些条件唯一确定（参见习题 10 d) 与 7 c)）。
\end{enumerate}

¶ 16) 假设第 3 小节的诸条件成立，此外 A = K(i)。设 u 是 E 的一个自同态。

\begin{enumerate}
\def\labelenumi{\alph{enumi})}
\item
  \(u\) 的特征值所成的集合含于 \(\Phi(x, u(x))\) 当 \(x\) 遍历 E 中满足 \(\Phi(x, x) = 1\) 的元素所成之集时所取值的集合 U。
\item
  我们称 A = K(i) 的一个子集 C 是凸的，如果对每对元素 \((\xi, \eta) \in \mathrm{C}^{2}\) 和每个 \(\tau \in K\)（满足 \(0 \leqslant \tau \leqslant 1\)），有 \(\tau\xi + (1 - \tau)\eta \in C\)。证明集合 U 是凸的。（化归到 n = 2 的情形；把 u 写成 \(\nu + iw\) 的形式，其中 \(\nu\) 与 w 是埃尔米特的，并在必要时把 u 换成 \(\lambda u\)（其中 \(\lambda \in A\)，\(\lambda\overline{\lambda} = 1\)），证明一切都化归为证明下列性质。设 \(f(\xi_{1}, \xi_{2}) = \xi_{1}\overline{\xi}_{1} + \xi_{2}\overline{\xi}_{2}\)、\(g(\xi_{1}, \xi_{2}) = a\xi_{1}\overline{\xi}_{1} + b\xi_{2}\overline{\xi}_{2}\)（\(a \in K, b \in K\)）、\(h(\xi_{1}, \xi_{2}) = \alpha\xi_{1}\overline{\xi}_{1} + \beta\xi_{1}\overline{\xi}_{2} + \beta\xi_{2}\overline{\xi}_{1} + \gamma\xi_{2}\overline{\xi}_{2}\)（\(\alpha \in K, \gamma \in K, \beta \in A\)）；设 \((\eta_{1}, \eta_{2}) \in A^{2}\)、\((\zeta_{1}, \zeta_{2}) \in A^{2}\) 满足 \(f(\eta_{1}, \eta_{2}) = f(\zeta_{1}, \zeta_{2}) = 1\)、\(g(\eta_{1}, \eta_{2}) = g(\zeta_{1}, \zeta_{2}) = 1\)、\(h(\eta_{1}, \eta_{2}) > 0\)、\(h(\zeta_{1}, \zeta_{2}) < 0\)；则存在 \((\theta_{1}, \theta_{2}) \in A^{2}\) 满足 \(f(\theta_{1}, \theta_{2}) = 1\)、\(g(\theta_{1}, \theta_{2}) = 1\) 与 \(h(\theta_{1}, \theta_{2}) = 0\)。可以先指出，对每对 \((\xi_{1}, \xi_{2})\)，存在 \(\mu \in A\) 满足 \(\mu\overline{\mu} = 1\) 与 \(\beta\mu\xi_{1}\overline{\xi}_{2} + \beta\overline{\mu}\xi_{2}\overline{\xi}_{1} = 0\)，这将使我们可以化归到 \(\beta = 0\) 的情形；利用第六章 § 2 第 5 小节命题 5。）
\item
  证明若 \(u\) 是正规的，则 U 是包含 \(u\) 的所有特征值的最小凸集。给出一个 \(u\) 不是正规的但 U 仍具有上述性质的例子。（取 \(u\) 的特征值为元素 \(\pm 1 \pm i\)，它们是单特征值，并取 0 作为一个二重特征值。）
\end{enumerate}

¶ 17) 我们假设第 3 小节的诸条件成立；对每个 \(\xi\in\mathbf{A}\)，我们用 \(|\xi|\) 记 K 的元素 \(\rho\geqslant0\)，它满足 \(\rho^{2}=\xi\bar{\xi}\)（\(\xi\) 的绝对值）。对 A 上每个方阵 \(M=(\alpha_{ij})\)，我们置 \(f(M)=\max_{i}|\alpha_{ii}|\)、\(g(M)=\max_{i,j}|\alpha_{ij}|\)，并用 \(\varphi(M)\) 记 \(M\) 的特征值的最大绝对值（将证明当 A 是 K 上的四元数体时这个定义是有意义的）。

\begin{enumerate}
\def\labelenumi{\alph{enumi})}
\tightlist
\item
  设 \(A, B, D\) 是 A 上三个 \(n\) 阶方阵。证明：若 \(D\) 是对角矩阵，则
\end{enumerate}

\[
g ^ {2} (A D B ^ {*}) \leqslant f ^ {2} (D) f (A A ^ {*}) f (B B ^ {*})\tag{1}
\]

（利用第 1 小节的不等式 (1)）。由此推出

\[
g (A B B ^ {*} A ^ {*}) \leqslant f (A A ^ {*}) \varphi (B B ^ {*})\tag{2}
\]

（把命题 5 应用于埃尔米特矩阵 \(BB^{*}\)）。由此推出：对 A 上 \(m\) 个任意方阵 \(A_{i}\)（\(1 \leqslant i \leqslant m\)），我们有

\begin{enumerate}
\def\labelenumi{(\arabic{enumi})}
\setcounter{enumi}{2}
\tightlist
\item
\end{enumerate}

\[
g ^ {2} (A _ {1} A _ {2} \dots A _ {m}) \leqslant f (A _ {1} A _ {1} ^ {*}) \varphi (A _ {2} A _ {2} ^ {*}) \dots \varphi (A _ {m - 1} A _ {m - 1} ^ {*}) f (A _ {m} ^ {*} A _ {m})\tag{4}
\]

\[
\varphi^ {2} (A _ {1} A _ {2} \dots A _ {m}) \leqslant \varphi (A _ {1} A _ {1} ^ {*}) \varphi (A _ {2} A _ {2} ^ {*}) \dots \varphi (A _ {m} A _ {m} ^ {*})\tag{5}
\]

\[
g ^ {2} (A _ {1} A _ {2} \dots A _ {m}) \leqslant \varphi (A _ {1} A _ {1} ^ {*}) \dots \varphi (A _ {m} A _ {m} ^ {*}).
\]

（对 (3)，从 (2) 出发用归纳法论证。对 (4)，利用习题 12；最后由 (3) 推出 (5)。）

  证明当把 \(A_{m}^{*}A_{m}\) 换成 \(A_{m}A_{m}^{*}\) 时（注意可以有 \(f(A^{*}A)\neq f(A A^{*})\)），或者当把 \(\varphi(A_{i}A_{i}^{*})\) 换成 \(f(A_{i}A_{i}^{*})\)（对 \(2\leqslant i\leqslant m-1\)）时（取所有 \(A_{i}\) 都等于全元素为 1 的方阵），不等式 (3) 不再必然成立。

\begin{enumerate}
\def\labelenumi{\alph{enumi})}
\setcounter{enumi}{1}
\tightlist
\item
  设 u 是 E 的一个正规自同态；若 λ 是 u 的一个特征值（当 A = K 时，是 K(i) 的一个元素），则 \(\lambda\bar{\lambda}\) 是 u*u 的一个特征值。因此，对每个正规矩阵 M（E 的一个正规自同态关于一个标准正交基的矩阵），我们有 φ²(M) = φ(MM*)。由此推出也有 g(M) ≤ φ(M)，且更一般地，若 M₁, . . ., Mₘ 是正规的，则
\end{enumerate}

\begin{enumerate}
\def\labelenumi{(\arabic{enumi})}
\setcounter{enumi}{5}
\tightlist
\item
\end{enumerate}

\[
g (M _ {1} \dots M _ {m}) \leqslant \varphi (M _ {1}) \dots \varphi (M _ {m})\tag{7}
\]

\[
\varphi (M _ {1} \dots M _ {m}) \leqslant \varphi (M _ {1}) \dots \varphi (M _ {m})
\]

（利用 (4) 与 (5)）。

\begin{enumerate}
\def\labelenumi{\alph{enumi})}
\setcounter{enumi}{2}
\tightlist
\item
  证明若 \(H\) 是 A 上一个正埃尔米特矩阵，我们有 \(f(H) = g(H) \leqslant \varphi(H)\)（利用 (3) 与 b)，把 H 写成 \(H = AA^{*}\)）。
\end{enumerate}

¶ 18) 我们假设第 3 小节的诸条件成立。

\begin{enumerate}
\def\labelenumi{\alph{enumi})}
\tightlist
\item
  对 E 的每个自同态 u，设 \((e_{i})\) 是由 \(u^{*}u\) 的特征向量组成的一个 E 的标准正交基，\(\rho_{i}\) 是 \(u^{*}u\) 对应于向量 \(e_{i}\) 的特征值；我们置 \(s(u) = (\sum_{i=1}^{n} \rho_{i})^{1/2}\)（当 A 交换时，这是 \(\mathrm{Tr}(u^{*}u)\) 的平方根）；我们有 \(s(u^{*}) = s(u)\)。若 u、v 是 E 的两个自同态，U、V 是它们关于 E 的同一标准正交基的矩阵，证明对元素在 K 中的矩阵 \(UV^{*} + VU^{*}\)，有 \((\mathrm{Tr}(UV^{*} + VU^{*}))^{2} \leqslant 4s(u)s(v)\)（注意 \((u^{*} + \lambda v^{*})(u + \lambda v)\) 对每个 \(\lambda \in K\) 是正埃尔米特的）；由此推出 \(s(u + v) \leqslant s(u) + s(v)\)。若此外 A 是交换的，则有
\end{enumerate}

\[
\left| \operatorname{Tr} (u v) \right| ^ {2} \leqslant s (u) s (v) \quad \text { and } \quad \left| \operatorname{Tr} (u) \right| \leqslant \sqrt {n} s (u)
\]

（利用习题 14 b) 把 \(u\) 写成一个乘积）。

\begin{enumerate}
\def\labelenumi{\alph{enumi})}
\setcounter{enumi}{1}
\tightlist
\item
  我们进一步假设 A 是交换的（故等于 K 或 K(i)）。若 \(H_{1}, H_{2}\) 是两个正埃尔米特方阵，证明有 \(\mathrm{Tr}(H_{1}H_{2}) \geqslant 0\)（化归到 \(H_{2}\) 是对角矩阵的情形）。由此推出我们有（用习题 17 的记号）
\end{enumerate}

\[
\left| \operatorname{Tr} (H _ {1} H _ {2}) \right| \leqslant \varphi (H _ {1}) \operatorname{Tr} (H _ {2}) \leqslant \operatorname{Tr} (H _ {1}) \operatorname{Tr} (H _ {2}).
\]

  从这些不等式推出：对 E 的任意两个自同态 \(u, \nu\)，我们有 \(s(u\nu) \leqslant s(u)s(\nu)\)。

\begin{enumerate}
\def\labelenumi{\alph{enumi})}
\setcounter{enumi}{2}
\tightlist
\item
  仍假设 A 是交换的，设 \(\prod_{i=1}^{n}(x-\lambda_i)\) 是 E 的任一自同态 \(u\) 的特征多项式的分解为线性因式之积；证明有 \(\sum_{i=1}^{n}|\lambda_i|^2\leqslant(s(u))^2\)，且要使不等式两边相等，必须且只需 \(u\) 是正规的（利用习题 12）。
\end{enumerate}

\begin{enumerate}
\def\labelenumi{\arabic{enumi})}
\setcounter{enumi}{18}
\item
  假设第 3 小节的诸条件成立。设 M 是 A 上一个 n 阶方阵；证明对 M 的每个方子矩阵 N（从 M 中删去若干行以及与这些行同指标的列而得到），有（用习题 17 的记号）\(\varphi^{2}(N) \leqslant \varphi(MM^{*})\)（适当地应用习题 17 的公式 (4)）。特别地，若 M 是正规矩阵，则 \(\varphi(N) \leqslant \varphi(M)\)（参见习题 17 b)）。
\item
  我们假设第 3 小节的诸条件成立，此外 A 等于 K 或 K(i)。把习题 17 至 19 的结果应用于 Φ 到 p 次外幂的扩张（§ 1 第 9 小节）以及所考虑的自同态或矩阵的 p 次外幂。特别地，证明：若 \(\prod_{i=1}^{n}(X-\lambda_i)\) 与 \(\prod_{i=1}^{n}(X-\rho_i^2)\) 分别是某个
\end{enumerate}

E 的自同态 \(u\) 与自同态 \(u^*u\) 的特征多项式的线性因式分解，且若假设 \(|\lambda_i| \geqslant |\lambda_{i+1}|\) 与 \(\rho_i \geqslant \rho_{i+1} \geqslant 0\)（\(1 \leqslant i \leqslant n-1\)），则有

\[
\mid \lambda_ {1} \lambda_ {2} \dots \lambda_ {h} \mid \leqslant \rho_ {1} \rho_ {2} \dots \rho_ {h}
\]

对 \(1 \leqslant h \leqslant n - 1\) 成立，且有 \(|\lambda_1\lambda_2 \ldots \lambda_n| = \rho_1\rho_2 \ldots \rho_n\)。

\begin{enumerate}
\def\labelenumi{\arabic{enumi})}
\setcounter{enumi}{20}
\tightlist
\item
  我们假设第 3 小节的诸条件成立。设 u 是 E 的一个埃尔米特自同态，\(\prod_{i=1}^{n}(X-\lambda_{i})\) 是它的特征多项式的线性因式分解；我们假设 \(\lambda_{i}\geqslant\lambda_{i+1}\)（\(1\leqslant i\leqslant n-1\)）。
\end{enumerate}

\begin{enumerate}
\def\labelenumi{\alph{enumi})}
\item
  证明 K 中的特征值 \(\lambda_{i}\) 中最大者（相应地，最小者）等于 \(\Phi(u(x), x)\) 的值中最大者（相应地，最小者），这里 \(x\) 遍历由这样的 \(x \in E\) 所成的集合：\(\Phi(x, x) = 1\)。（直接论证，或者化归到 A = K(i) 的情形后应用习题 16 c)（ \(\S\) 3 习题 4）。）
\item
  设 \(\Psi_{\mathrm{v}}\) 是埃尔米特型 \(\Psi\)（与 \(u\) 相伴的）在 E 的一个向量子空间 V 上的限制，\(u_{\mathrm{v}}\) 是与 \(\Psi_{\mathrm{v}}\) 相伴的 V 的埃尔米特自同态。证明 \(\lambda_{k}\) 是诸 \(u_{\mathrm{v}}\)（当 V 遍历 E 的维数为 \(n - k + 1\) 的向量子空间所成之集时）的最大特征值中最小的一个（利用命题 5）。
\end{enumerate}

\begin{enumerate}
\def\labelenumi{\arabic{enumi})}
\setcounter{enumi}{21}
\tightlist
\item
  我们假设第 3 小节的诸条件成立，此外 A 等于 K(i) 或 K 上的四元数体。设 u、v 是 E 的两个正规自同态；设 (E \(_{i}\) ) \(_{1\leqslant i\leqslant r}\)（相应地，(F \(_{j}\) ) \(_{1\leqslant j\leqslant s}\)）是 E 分解为两两正交的子空间之直和，使得在 E \(_{i}\)（相应地，F \(_{j}\)）中存在一个由 u（相应地，v）的属于同一特征值 \(\lambda_{i}\)（相应地，\(\mu_{j}\)）的特征向量组成的标准正交基，且 \(\lambda_{h}\) 与 \(\lambda_{i}\)（相应地，\(\mu_{j}\) 与 \(\mu_{k}\)）当 \(h\neq i\)（相应地，\(j\neq k\)）时不能通过 A 的一个内自同构互相变换（参见命题 4 与习题 9）。要使 E 的一个自同态 w 满足 uw = wν，必须且只需对每个 j（\(1\leqslant j\leqslant s\)），w(F \(_{j}\) ) 含于某个 E \(_{i}\) 之中，且 w 使 v 的属于特征值 \(\mu_{j}\) 的每个特征向量的像是 u 的属于特征值 \(\mu_{j}\) 的一个特征向量（特别地，这蕴涵 \(\mu_{j}\) 与 \(\lambda_{i}\) 可通过 A 的一个内自同构互相变换）。由此推出：若确实如此，则有 u*w = wν*。
\end{enumerate}

¶ 23) 我们假设第 3 小节的诸条件成立。设 u 与 v 是 E 的两个自同态。

\begin{enumerate}
\def\labelenumi{\alph{enumi})}
\item
  我们假设 u 与 uv 是正规的。要使 vu 是正规的，必须且只需 v 与 \(u^{*}u\) 交换。（要看出条件是必要的，利用关系式 \(u(vu) = (uv)u\) 与习题 22；要看出条件是充分的，利用习题 14 a) 与 10 d)。）
\item
  我们假设 \(u, v\) 与 \(uv\) 正规；设 \(\beta\) 是 \(uu^*\) 的最大特征值，\(F\) 是 \(E\) 的由 \(uu^*\) 的属于特征值 \(\beta\) 的特征向量所成的子空间。证明 \(v(F) \subset F\)。（注意对每个正规自同态 \(u\) 和每个 \(x \in E\)，有
\end{enumerate}

\[
\Phi (u (x), u (x)) = \Phi (u ^ {*} (x), u ^ {*} (x)).
\]

  化归到 A 等于 K(i) 或 K 上的四元数体的情形；于是 F 具有一个由 u（与 \(u^{*}\)）的特征向量组成的基；注意对这样的向量 z，我们有 \(\Phi(u^{*}u\nu(z),\nu(z))=\beta\Phi(\nu(z),\nu(z))\)。由此推出 \(\nu u\) 是正规的（对 \(u^{*}u\) 的不同特征值的个数作归纳论证）。若 h（相应地，\(h'\)）是使 \(h^{2}=uu^{*}\)（相应地，\(h'^{2}=\nu\nu^{*}\)）的正埃尔米特自同态，且若置 \(u=hu_{1}\)、\(\nu=h'\nu_{1}\)，其中 \(u_{1}\) 与 \(\nu_{1}\) 是酉的，h 与 \(u_{1}\) 交换且 \(h'\) 与 \(\nu_{1}\) 交换（习题 14 a)），证明偶 \((h,h')\)、\((h,\nu_{1})\) 与 \((h',u_{1})\) 都是可交换的；反之亦然。由此推出 \(u^{m}\nu^{n}\)、\(\nu^{n}u^{m}\)、\(u\nu^{*}\) 与 \(\nu^{*}u\) 这时都是正规的（m 与 n 是任意整数 \textgreater0）。

¶ 24) 假设第 3 小节的诸条件成立。设 Γ 是 E 的一个自同构群，使得每个 \(u \in \Gamma\) 都是正规的。证明存在 E 的一个分解，分解为两两正交的子空间 \(E_k\)（\(1 \leqslant k \leqslant r\)）的直和，使得 \(E_k\) 上每个 \(u \in \Gamma\) 的限制都具有 \(\lambda_k v_k\) 的形式，其中 \(\lambda_k\) 是 K 的一个元素 \textgreater0，\(v_k\) 是 \(E_k\) 的一个酉自同态（把每个 \(u \in \Gamma\) 分解成 \(h\nu\) 的形式，其中 \(v\) 是酉的，\(h\) 是正埃尔米特的，且 \(h^2 = uu^*\)（习题 14a)；利用习题 23b 并把定理 2 的推论 2 应用于由埃尔米特自同态 \(h\)（它们对应于诸 \(u \in \Gamma\)）生成的代数）。由此推出：若 Γ 是有限的，则诸元素 \(u \in \Gamma\) 都是酉的。

\begin{enumerate}
\def\labelenumi{\arabic{enumi})}
\setcounter{enumi}{24}
\tightlist
\item
  假设 A = K（极大有序域）。设 L 是 A 上一个 n 维欧几里得空间，其度量型是正且非退化的。
\end{enumerate}

  设 S 是 L 中一个非退化仿射二次曲面（§ 6 习题 25）。若 S 具有一个中心 a，且在 L 中取 a 为原点，证明存在 L 的一个标准正交基 \((e_{i})\)，使得关于这个基，S 有方程 \(\lambda_{1}\xi_{1}^{2} + \cdots + \lambda_{n}\xi_{n}^{2} = 1\)。此外，若两个标准正交基都具有这个性质，则与它们对应的元素 \(\lambda_{i} \in K\) 除次序外是相同的。

  若 S 没有中心，证明存在 S 的一个点 b 以及（取 b 为原点）L 的一个标准正交基，使得关于这个基，S 有方程 \(\lambda_{1}\xi_{1}^{2} + \cdots + \lambda_{n-1}\xi_{n-1}^{2} + \xi_{n} = 0\)。（若 \(\overline{S}\) 是使 \(S = L \cap \overline{S}\) 的射影二次曲面，c 是关于 \(\overline{S}\) 的无穷远超平面 \(H_{0}\) 的（无穷远）极点，则由下列条件确定 b：过 b 与 c 的直线垂直于 S 在点 b 处的切超平面。）

¶ 26) a) 设 K 是一个交换域，S 是 K 的一个子集，使得 K 是极大 S 可序域（第六章 § 2 习题 8）。设 f 是 K{[}X{]} 中一个多项式，L 是它的分裂域。证明：若对 K 上每个与其环结构相容的（全）序结构，L 都容许 K 的一个有序扩张结构，则必有 L = K。（用反证法论证：按第六章 § 2 习题 8 e) 的方式进行，化归到将有 {[}L : K{]} = 2 的情形。最后注意，若 b ∈ K 不是平方，则存在 K 的一个与其环结构相容的全序结构，使 b \textless{} 0（参见第六章 § 2 第 3 小节定理 1 的引理）。）

\begin{enumerate}
\def\labelenumi{\alph{enumi})}
\setcounter{enumi}{1}
\tightlist
\item
  把第 3 小节的结果（命题 6 除外）推广到 K 是极大 S 可序域的情形 (*)。（先利用 a) 证明命题 5 的类似命题。然后对交换代数 B 由埃尔米特自同态组成的情形，利用第八章 § 9 第 4 小节命题 10，建立定理 2 推论 2 的类似命题。对 A = K(i) 时正规自同态 u 的情形，注意这时可以把 u 写成 u = v + iw，其中 v 与 w 是埃尔米特的且交换。最后，若 A 是 K 上的四元数体，E₀ 是把 E 看作 K(i) 上维数 2n 的向量空间所得的空间，且若置 Φ(x, y) = Φ₁(x, y) + Φ₂(x, y)j，其中 Φ₁ 与 Φ₂ 取值在 K(i) 中，则注意若 u 对 Φ 是正规的，它对 Φ₁ 也是正规的。）
\end{enumerate}

¶ 27) 设 K 是一个极大有序域，E 是 K 上一个 n 维向量空间，Φ 是 E 上一个非退化对称双线性型，其符号差 (s, t) 不同于 (n, 0) 与 (0, n)。设 (ei) 是 Φ 的一个正交基，使得对 1 ≤ i ≤ s 有 Φ(ei, ei) = 1，对 s + 1 ≤ i ≤ n 有 Φ(ei, ei) = -1。对每个正交变换 u ∈ U(Φ)，设

\[
U = \left( \begin{array}{c c} M & N \\ P & Q \end{array} \right)
\]

\(u\) 关于 \((e_{i})\) 的矩阵，把它写成矩阵的方块阵列的形式，对应于把 \([1, n]\) 分为 \([1, s]\) 与 \([s + 1, n]\) 的分划。

\begin{enumerate}
\def\labelenumi{\alph{enumi})}
\tightlist
\item
  证明下列关系式
\end{enumerate}

\[
\begin{array}{r l} ^ {t} M. M - ^ {t} P. P & = I _ {s} \\ ^ {t} Q. Q - ^ {t} N. N & = I _ {t} \\ ^ {t} M. N - ^ {t} P. Q & = 0. \end{array}
\]

\begin{enumerate}
\def\labelenumi{\alph{enumi})}
\setcounter{enumi}{1}
\item
  设 R 是 K 上一个 t 行 s 列矩阵，使得 \(I_{s} - ^{t}R \cdot R\) 是一个正非退化埃尔米特型（在 \(K^{s}\) 上）的矩阵。证明 det \((M + \lambda NR)\) 在 K 中对 \(-1 \leqslant \lambda \leqslant 1\) 不变号（利用 a)，证明 \(^{t}(M + \lambda NR)(M + \lambda NR)\) 是一个正非退化对称型的矩阵）。
\item
  置 \(\sigma(u) = \sigma(U) = \operatorname{sgn}(\det M)\)；证明对任意两个元素 \(u\)、\(\nu\)（属于正交群 \(\mathbf{U}(\Phi)\)），有 \(\sigma(u\nu) = \sigma(u)\sigma(\nu)\)（利用 \(a\) 与 \(b\)）。由此推出特殊正交群 \(\mathbf{SU}(\Phi)\) 的换位子群不同于 \(\mathbf{SU}(\Phi)\)（参见 § 10 习题 9）。
\end{enumerate}

